%%%PAGE 77 rot=0 legibility=good summary=系统动能定理与机械能守恒公式；连接体能量方法；上抛下落类（有双向性）结论及斜面推广；小球竖直上抛 Ek=2Ep、Ep=2Ek 求 h 例题
%%%BLOCK id=077-01 ch=06 sec=动能定理与机械能守恒·系统 kind=formula alt=none cont=none y=0.07-0.20
\textbf{系统动能定理．}
\[
\sum W_{\text{外}}+\sum W_{\text{内}}=\Delta E_k
\qquad
\text{内力}
\begin{cases}
W\ \text{绳、杆、}f_{\text{静}}\ \text{抵消}\\
W\ \red{\underline{\text{弹性绳}}}\text{、}\red{\underline{\text{弹簧}}}\text{、}\red{\underline{\text{不抵消}}}
\end{cases}
\]
% 原稿：第二行“弹性绳”“弹簧”“不抵消”下为红笔下划线（字为黑色）
\textbf{机械能守恒}
\[
W_{\text{其}}+W_G+W_{\text{弹}}=\Delta E_k
\qquad
W_{\text{其}}=\Delta E_{k}+\Delta E_{pG}+\Delta E_{p\text{弹}}=\Delta E_{\text{机}}
\]
% CHECK: 第二式 ΔE_k 后手写多一小笔（像逗号或下标 1）
\[
W_{\text{其}}=0\ \text{时 机械能守恒}
\]
%%%END
%%%BLOCK id=077-02 ch=06 sec=连接体·能量 kind=method alt=03 cont=none y=0.19-0.31
\textbf{连接体}
\[
\begin{cases}
\text{对1 由动能定理}\\
\text{对2 由动能定理}\\
\text{速度关联}
\end{cases}
\qquad
\fcolorbox{red!75!black}{white}{\begin{tabular}{@{}c@{}}单体不守恒\\ 系统\,才守恒\end{tabular}}
\]
% 上面“单体不守恒/系统才守恒”外为红圈；下面三行外有红笔竖括号
\begin{itemize}
\item 过程为变加变减速，用超失重分析（对××力与 $mg$ 比较）
\item $V_1\ \overset{\text{\scriptsize 正功}}{0}\ \text{有}\ \overset{\text{\scriptsize 负功}}{0}$
\item $V_2\ \text{有}\ \overset{\text{\scriptsize 负功}}{0}\ \overset{\text{\scriptsize 正功}}{\text{有}}$
\end{itemize}
% 小字“正功/负功”写在对应符号上方；第一行正功在第一个0上、负功在最后的0上；第二行负功在“有”“0”之间偏0、正功在最后的“有”上
弹簧速度不能分解
%%%END
%%%BLOCK id=077-03 ch=06 sec=解题分类·过程与方程 kind=summary alt=none cont=none y=0.31-0.42
不以模型分过程而以方程分过程．
\[
\begin{cases}
\text{全过程类}\\
\text{上抛下落类}\\
\text{绳杆连接体类．}
\end{cases}
\]
%%%END
%%%BLOCK id=077-04 ch=06 sec=上抛下落类·双向性 kind=method alt=none cont=none y=0.45-0.62
\textbf{上抛下落类}\quad （有双向性）．
\begin{keybox}
结论\quad 无论 $f$ 是否存在
\end{keybox}
\red{能量 $\times n$\quad 速度 $\times\sqrt{n}$}
% 以下两行为原稿右侧的黑字小注
\[
E_{k0}\times n\quad h\times n\quad E_{k\text{返}}\times n\quad W_f\times n\quad W_G\times n
\]
\[
V_{n\text{返}}\times\sqrt{n}
\]
% CHECK: “返”字辨认存疑（E_k返、V_n返）
\begin{align*}
\text{上}\quad & 0-\tfrac12 mV_0^2=-mgh-fh \quad\text{①}\\
\text{下}\quad & \tfrac12 mV^2-0=mgh-fh \quad\text{②}\\
& \frac{\text{①}}{\text{②}}=\frac{V_0^2}{V^2}=\frac{f+mg}{mg-f}\\
& \tfrac12 mv^2-\tfrac12 mv_0^2=-2hf
\end{align*}
%%%END
%%%BLOCK id=077-05 ch=06 sec=上抛下落类·双向性 kind=method alt=02 cont=none y=0.62-0.80
\red{$0-V_{\text{上}}^2=-g\sin\theta-\mu g\cos\theta$}\\
\red{$V_{\text{下}}^2=g\sin\theta-\mu g\cos\theta$}
% CHECK: 红字 V 的下标手写像 E/F，按上下文取 上/下
%FIG p077_a 0.065 0.72 0.205 0.78
\notefig[0.25]{figures/p077_a.png}
\[
\begin{cases}
W_G = mg\sin\theta\,S\\
W_f = \mu mg\cos\theta\,S
\end{cases}
\qquad
\frac{W_G}{W_f}=\frac{\tan\theta}{\mu}\ \text{为定值}
\qquad
W_f\propto W_G
\]
%%%END
%%%BLOCK id=077-06 ch=06 sec=上抛下落类·双向性 kind=problem alt=none cont=none y=0.81-0.99
\begin{problem}
\unclear{小}球竖直上抛，最高 $H$，$f$ 恒定，上升至 $h$ 处，$E_k=2E_p$，下落至离地 $h$ 处，$E_p=2E_k$，求 $h$。
\begin{solution}
%FIG p077_b 0.0 0.86 0.28 0.97
\notefig[0.3]{figures/p077_b.png}
\[
\left\{
\begin{array}{l}
\uparrow\ E_k=2mgh\qquad \downarrow\ E_k'=\tfrac12 mgh\\[6pt]
h\ \left\{
\begin{array}{l}
-\tfrac32 mgh=-W_f\times 2\\
W_f=-\tfrac34 mgh\\
0-2mgh=-\tfrac34 mgh-mg(H-h)\\
\therefore h=\tfrac49 H
\end{array}
\right.
\end{array}
\right.
\]
% 原稿左侧示意：顶部一条横线，↑自小圆圈向上、↓向下止于小圆圈（上升/下落经过 h 处）
% CHECK: −(3/2)mgh=−W_f×2 与 W_f=−(3/4)mgh 符号不一致（原稿如此）
\end{solution}
\end{problem}
%%%END
%%%PAGEEND 77
